Soit $\mathcal{S}$ un ensemble fini de points du plan, contenant au moins deux points. On suppose que trois
points quelconques de $\mathcal{S}$ ne sont pas alignés.

On appelle \emph{moulin à vent} le processus suivant : le processus commence avec une droite $\ell$ contenant
un unique point $P$ de $\mathcal{S}$ ; la droite $\ell$ tourne, dans le sens des aiguilles d'une montre, autour
du point $P$, appelé \emph{pivot}, jusqu'à ce qu'elle rencontre pour la première fois un autre point de
$\mathcal{S}$ ; ce point, $Q$, devient le nouveau pivot ; la droite continue alors sa rotation dans le sens des
aiguilles d'une montre autour de $Q$, jusqu'à rencontrer un nouveau point de $\mathcal{S}$ ; ce processus
continue indéfiniment.

Montrer qu'on peut choisir un point $P$ de $\mathcal{S}$ et une droite $\ell$ contenant $P$, de façon que le
moulin à vent commençant par $\ell$ utilise chaque point de $\mathcal{S}$ comme pivot une infinité de fois.
